\section{Subjective probability and dynamics: finite derivations}
\label{sec:p15-rat41-annex}

These proofs fix finite scopes for mathematical claims in the selected
H\'ajek--Staffel survey \parencite[pp.~239--248]{HajekStaffel2021Dynamics}.
They do not prove the normative authority of the surveyed principles.

\begin{annexitem}{annex:p15.rat41.brier}{thm:p15-rat41-brier}{Properness, finite dominance, and prior-optimal planning}
\begin{proof}
Expanding the binary expectation gives
\(q(1-x)^2+(1-q)x^2=q-2qx+x^2=q(1-q)+(x-q)^2\).
The last square vanishes exactly at \(x=q\). Absolute error gives
\(q(1-x)+(1-q)x=q+(1-2q)x\), proving the stated endpoint and tie cases.

The set \(K\) is nonempty and compact: it is the image of the finite
probability simplex under the continuous averaging map. It is convex.
The continuous squared-distance function therefore has a minimizer
\(y\in K\). Write \(h=c-y\). For any \(z\in K\), the segment
\(y+u(z-y)\) remains in \(K\), \(0\leq u\leq1\). The right derivative
at zero of its squared distance from \(c\) is
\(-2h\cdot(z-y)\), which cannot be negative. Thus
\(h\cdot(z-y)\leq0\). For every truth vector \(v\in K\),
\[
 \|c-v\|^2-\|y-v\|^2
       =\|h\|^2+2h\cdot(y-v)\geq\|h\|^2>0
\]
when \(c\notin K\). This proves strict worldwise Brier improvement.
The minimizer is unique: if two distinct minimizers existed, their
midpoint's squared distance would equal the average of their squared
distances minus a quarter of their squared mutual distance, contradicting
minimality.

If \(c\in K\), choose probability masses on worlds whose truth-vector
average is \(c\). For any proposed replacement \(b\), averaging its
score difference over these masses gives
\[
                       \mathbb E[S(b,W)-S(c,W)]=\|b-c\|^2\geq0.
\]
A strictly smaller score at every world would have a strictly smaller
average because the finite probability masses sum to one. Hence no
replacement strictly dominates \(c\).

Finally fix the prior and observation partition. Conditional on a cell
of positive prior mass, each event has the old conditional probability
\(q\); the first identity makes its unique optimal report \(q\).
Sum over events and average over cells. Every positive-weight squared
deviation adds nonnegative expected loss, and adds strictly positive loss
when any report differs. The old conditional reports form a coherent
conditional probability and so remain feasible even when planned reports
are required to be coherent. Null cells have zero prior weight and impose
no expected-score constraint.
\end{proof}
\end{annexitem}

\begin{annexitem}{annex:p15.rat41.positive-atoms}{thm:p15-rat41-positive-atoms}{The countable positive-singleton bound}
\begin{proof}
For each integer \(k\geq1\), at most \(k\) disjoint singletons can
have mass at least \(1/k\): otherwise finite additivity and nonnegativity
would give a finite union mass greater than one. Every singleton of
positive mass belongs to one such finite set, by choosing \(k\) large
enough that \(1/k\) does not exceed its mass. The positive-mass
singletons therefore form a countable union of finite sets. An
uncountable measurable singleton family cannot be entirely contained
in this countable set.
\end{proof}
\end{annexitem}

\begin{annexitem}{annex:p15.rat41.indifference}{thm:p15-rat41-indifference}{Entropy and minimax properties on a fixed finite partition}
\begin{proof}
First recall finite Gibbs nonnegativity. For probability vectors \(a,b\),
if \(a_j>0=b_j\) for some \(j\), the relative entropy is infinite.
Otherwise \(-\ln u\geq1-u\) for \(u>0\), with equality only at one,
gives
\[
 \sum_{a_j>0}a_j\ln(a_j/b_j)
 \geq\sum_{a_j>0}(a_j-b_j)
 =1-\sum_{a_j>0}b_j\geq0.
\]
Equality requires \(a_j=b_j\) on positive \(a\) coordinates and no
remaining mass in \(b\), hence \(a=b\). Divide by \(\ln2\) to
obtain the base-two result.

For the uniform vector \(u_j=1/n\), the divergence is
\(D(p,u)=\log_2n-H(p)\). Gibbs nonnegativity and its equality
condition prove the entropy maximum and its uniqueness, and negation
gives the source's information minimum.

In outcome \(j\), the Brier score is
\(1-2p_j+\sum_i p_i^2\). Its maximum is at least its average over
the \(n\) outcomes, namely \(1-2/n+\sum_i p_i^2\).
Since \(\sum_i(p_i-1/n)^2=\sum_i p_i^2-1/n\geq0\), this is
at least \(1-1/n\), with equality only at the uniform vector.
That vector has the same score \(1-1/n\) in every outcome and attains
the lower bound. The argument includes \(n=1\).
\end{proof}
\end{annexitem}

\begin{annexitem}{annex:p15.rat41.jeffrey}{thm:p15-rat41-jeffrey}{Partition weighting and the explicit noncommuting resets}
\begin{proof}
For \(w\in E_i\), put
\(q_w=r_i C_{\rm old}(\{w\})/C_{\rm old}(E_i)\).
The masses are nonnegative; their sum on \(E_i\) is \(r_i\), and
their total is one. Summing over \(A\) gives the Jeffrey formula.
When \(r_i>0\), dividing the masses on \(E_i\) by \(r_i\) gives
exactly the old conditional law. Conversely, any probability with
these cell masses and conditional laws has these same point masses;
a zero-mass cell must have all its point masses zero. This proves
uniqueness and the certainty specialization.

In the example, both initial row and column marginals are \(1/2\).
Resetting the first row to \(4/5\) multiplies its entries by \(8/5\)
and the second row by \(2/5\), yielding
\[
                 \begin{pmatrix}16/25&4/25\\1/25&4/25\end{pmatrix}.
\]
Its first column mass is \(17/25\) and second column mass \(8/25\).
Resetting the first column to \(4/5\) therefore multiplies column
entries by \(20/17\) and \(5/8\), yielding
\(\bigl(\begin{smallmatrix}64/85&1/10\\4/85&1/10\end{smallmatrix}\bigr)\).
Reversing the operations transposes this final matrix because the
initial matrix is symmetric. The off-diagonal entries \(4/85\)
and \(1/10\) differ, so the final distributions differ. The first
sequence's first row sum is \(64/85+1/10=29/34\).
\end{proof}
\end{annexitem}

\begin{annexitem}{annex:p15.rat41.minimal-update}{thm:p15-rat41-minimal-update}{Oriented relative entropy, Hellinger, and variation}
\begin{proof}
Write \(c_i=P(E_i)>0\). Any \(Q\) assigning positive mass where
\(P\) is zero has infinite \(D(Q,P)\) and cannot minimize, since
the Jeffrey distribution has finite divergence. Otherwise, on a cell
with \(r_i>0\), factor its masses as \(q_w=r_i q_{w\mid i}\)
and \(p_w=c_i p_{w\mid i}\). Summing logarithms gives
\[
 D(Q,P)=\sum_{i:r_i>0}r_i\log_2(r_i/c_i)
       +\sum_{i:r_i>0}r_i D(Q(\cdot\mid E_i),P(\cdot\mid E_i)).
\]
The first term is fixed. Gibbs nonnegativity, proved in
Annex~\ref{annex:p15.rat41.indifference}, makes the second nonnegative, with equality exactly when each
positive-weight conditional law is retained. Zero-weight cells have
zero point masses. The Jeffrey law is therefore the unique minimizer.

The squared Hellinger expression equals
\(2-2\sum_w\sqrt{q_wp_w}\). On each cell Cauchy--Schwarz gives
\(\sum_{w\in E_i}\sqrt{q_wp_w}\leq\sqrt{r_i c_i}\).
For \(r_i>0\), equality holds exactly when the square-root vectors
are proportional, which yields \(q_w=r_i p_w/c_i\). For \(r_i=0\)
the entire cell is zero. Thus the affinity is uniquely maximized by
the Jeffrey law, proving unique Hellinger minimization.

For variation, the triangle inequality on every cell gives
\[
 \sum_{w\in E_i}|q_w-p_w|
        \geq\left|\sum_{w\in E_i}(q_w-p_w)\right|=|r_i-c_i|.
\]
In the Jeffrey law all changes within a cell have the same sign,
so equality is attained. Uniqueness does not follow. In the stated
three-world example the two proposed laws have absolute changes
respectively \((1/4,1/4,1/2)\) and \((1/2,0,1/2)\). Both totals
are one, hence both variation distances are \(1/2\), equal to the
cellwise lower bound.
\end{proof}
\end{annexitem}

\begin{annexitem}{annex:p15.rat41.graded}{thm:p15-rat41-graded}{Finite projection benefits with normalized stakes and exact repairs}
\begin{proof}
Use the unique projection \(y\) and \(h=c-y\) from the Brier proof.
The inequality \(h\cdot(z-y)\leq0\) for all \(z\in K\) implies
that \(y\) is also the projection of \(c_s=y+(1-s)h\):
\[
 \|c_s-z\|^2
 =(1-s)^2\|h\|^2+\|z-y\|^2-2(1-s)h\cdot(z-y)
 \geq(1-s)^2\|h\|^2.
\]
Thus its distance to \(K\) is \((1-s)\|h\|\). For every truth
vector \(v\), subtracting scores yields
\[
 S(c,v)-S(c_s,v)
   =s(2-s)\|h\|^2+2s h\cdot(y-v)>0.
\]
The second term is nonnegative and the first strictly positive.

At quote vector \(c\), a stake vector \(b\) has agent net payoff
\(b\cdot(v-c)\). Restrict \(\|b\|\leq1\); its guaranteed loss
is \(\min_v b\cdot(c-v)\). Since \(y\) is an average of truth
vectors, this minimum is no greater than
\(b\cdot(c-y)\leq\|c-y\|\). Choosing \(b=h/\|h\|\) gives
\[
 b\cdot(c-v)=\|h\|+h\cdot(y-v)/\|h\|\geq\|h\|.
\]
The maximum guaranteed loss is therefore precisely the distance.
At a coherent vector the probability-weighted expected payoff is
zero, so no positive guaranteed loss is possible; \(b=0\) attains
zero. Apply the same argument to \(c_s\), including \(s=1\).

For the complement pair, coherent vectors are \((a,1-a)\),
\(0\leq a\leq1\). Squared distance from \((0.3,0.4)\) is
\[
       (a-0.3)^2+(0.6-a)^2=2(a-0.45)^2+0.045.
\]
This proves the projection and distance. Absolute distance is
\(|a-0.3|+|0.6-a|\), minimized at its constant lower bound
\(0.3\) exactly when \(0.3\leq a\leq0.6\).

Finally let \(c=(0.3,0.4)\), \(s_0=0.7\), and
\(\bar c=c/s_0=(3/7,4/7)\). For any probability pair \(q\),
\[
 \sum_i q_i\log_2(q_i/c_i)=D(q,\bar c)-\log_2s_0.
\]
Gibbs nonnegativity makes its unique minimizer \(q=\bar c\).
The reference vector \(c\) itself is not normalized, which is why
this objective was explicitly labeled an extension rather than a
relative entropy between probability distributions.

For the reversed orientation, let \(q\) be a probability pair. If some
\(q_i=0\), the term \(c_i\log_2(c_i/q_i)\) is \(+\infty\) because
\(c_i>0\). Otherwise, since \(c=s_0\bar c\) and \(\sum_i\bar c_i=1\),
\[
 \sum_i c_i\log_2\frac{c_i}{q_i}
   =s_0\sum_i\bar c_i\Bigl(\log_2\frac{\bar c_i}{q_i}+\log_2s_0\Bigr)
   =s_0\bigl(D(\bar c,q)+\log_2s_0\bigr).
\]
As \(s_0>0\), Gibbs nonnegativity again makes \(q=\bar c\) the unique
minimizer. For every probability pair \(q\),
\(\sum_i(c_i-q_i)=s_0-1\) is constant, so adding it to the first
objective leaves its unique minimizer \(\bar c\) unchanged. Numerically
\(\bar c=(0.42857\ldots,0.57142\ldots)\), which rounds to
\((0.429,0.571)\).
\end{proof}
\end{annexitem}
